\begin{mistake}{2026-05-13}{01}

\begin{problemcontent}
设 \(f(x) = \cos|x| + x^2|x|\)，在 \(x = 0\) 处存在的最高阶导数的阶数为\underline{\hspace{2em}}。
\end{problemcontent}

\begin{solutioncontent}
分别判定各项在 \(x=0\) 处的可导阶数，取"短板"。

\textbf{第一项：}\(\cos|x| = \cos x\)（余弦为偶函数），无穷阶可导，不构成瓶颈。

\textbf{第二项：}\(x^2|x|\)，利用结论：\(x^n|x|\) 在 \(x=0\) 处恰好 \(n\) 阶可导，\(n+1\) 阶不存在。

理由：\(x^n|x|\) 分段为 \(\pm x^{n+1}\)，求 \(n\) 阶导后得 \(\pm(n+1)!\,x\)（左右斜率相同，\(n\) 阶导存在）；求 \(n+1\) 阶导后得常数 \(\pm(n+1)!\)（左右不等，不存在）。

故 \(x^2|x|\) 在 \(x=0\) 处恰好 2 阶可导，3 阶不存在。

两项取短板，最高阶导数的阶数为 \(\boxed{2}\)。
\end{solutioncontent}

\mistakeknowledge{
\begin{itemize}
  \item \textbf{核心公式/定理}：\(x^n|x|\) 在 \(x=0\) 处恰好 \(n\) 阶可导，\(n+1\) 阶导数不存在
  \item \textbf{易错点}：忘记 \(\cos|x|=\cos x\)（偶函数性质）；多项相加时最高可导阶数取各项的最小值（短板原则）
  \item \textbf{关联知识}：\(x^n\) 在 \(x=0\) 处无穷阶可导；含绝对值函数可导性需分段后逐阶验证左右导数
\end{itemize}}

\end{mistake}
